\begin{table}[H]
    \centering
    \caption{Matriks Translasi Persentil Standar ke Skor Mentah per Subtes ($M = 500$, $SD$ Nasional Tidak Resmi per Subtes)}
    \label{tab:ranking_subtes}
    \begin{spacing}{1.0}
    \resizebox{\textwidth}{!}{%
    \begin{tabular}{lcccccccccc}
        \toprule
        Titik \textit{Benchmark} & Persentil (\%) & Peringkat Nasional & PU & PPU & PBM & PK & LBI Saintek & LBI Soshum & LBE & PM \\
        \midrule
        M - 3 SD   & 99,87 & 845.375 & 285,44 & 200,00 & 200,00 & 200,00 & 200,00 & 200,00 & 200,00 & 200,00 \\
        M - 2 SD   & 97,72 & 827.260 & 356,96 & 300,00 & 300,00 & 300,00 & 300,00 & 300,00 & 300,00 & 300,00 \\
        M - 1 SD   & 84,13 & 712.213 & 428,48 & 400,00 & 400,00 & 400,00 & 400,00 & 400,00 & 400,00 & 400,00 \\
        M          & 50,00 & 423.259 & 500,00 & 500,00 & 500,00 & 500,00 & 500,00 & 500,00 & 500,00 & 500,00 \\
        M + 1 SD   & 15,87 & 134.305 & 571,52 & 600,00 & 600,00 & 600,00 & 600,00 & 600,00 & 600,00 & 600,00 \\
        M + 2 SD   & 2,28 & 19.258 & 643,04 & 700,00 & 700,00 & 700,00 & 700,00 & 700,00 & 700,00 & 700,00 \\
        M + 3 SD   & 0,13 & 1.143 & 714,56 & 800,00 & 800,00 & 800,00 & 800,00 & 800,00 & 800,00 & 800,00 \\
        \bottomrule
    \end{tabular}%
    }
    \end{spacing}
    \vspace{2mm}
    \begin{flushleft}
    \footnotesize\textit{Catatan.} Persentil (\textit{top \%}) identik lintas subtes pada baris yang sama secara desain. Hanya skor mentah yang berbeda antarsubtes karena perbedaan $SD$ nasional. Peringkat nasional dihitung dari jumlah peserta hadir, $N_{\text{nasional}} = 846.518$.
    \end{flushleft}
\end{table}
